\section{Практические грабли}
\subsection{Простое приложение на pthread}

\begin{frame}
	\tableofcontents[currentsection,currentsubsection]
\end{frame}

\begin{frame}{Что такое pthread}
	\begin{itemize}
		\item Стандартный интерфейса функций для работы с потоками (POSIX Threads).
		\item Есть реализации под Windows, Linux и другие ОС.
		\item Стандарт при разработке программ на C.
		\item Имена функций и типов начинаются с \t{pthread\_}.
		\item В Linux можно получить справку, набрав \t{man <имя функции>} в консоли.
		\item Под остальными "--- то же самое, но в гугле.
	\end{itemize}
\end{frame}

\begin{frame}[fragile]{Пример кода}
\begin{minted}{cpp}
void* worker(void* arg) {
    printf("Hello from thread! arg=%d\n", *(int*)arg);
    *(int*)arg += 10;
    return arg;
}
int main(void) {
    pthread_t id;
    int data = 1234;
    assert(pthread_create(&id, NULL, worker, &data) == 0);
    void* retval;
    assert(pthread_join(id, &retval) == 0);
    assert(retval == &data);
    printf("data is %d\n", data);
    return 0;
}
\end{minted}
\end{frame}

\begin{frame}{Как живут потоки}
	\begin{itemize}
		\item При создании потока при помощи \t{pthread\_create} указывается функция и её аргумент "--- один указатель на что угодно.
		\item Вернуть функция тоже может указатель на что угодно.
		\item Поток завершается, когда функция делает \t{return} или \t{pthread\_exit}.
		\item Указатель на поток хранится в переменной типа \t{pthread\_t}.
		\item При создании потока он сразу начинает выполняться.
		\item \t{pthread\_join} делает следующее:
			\begin{enumerate}
				\item Ждёт окончания работы потока.
				\item Освобождает все ресурсы потока (стек).
				\item Возвращает то, что вернула функция потока.
			\end{enumerate}
		\item Когда \t{main} делает \t{return 0} или вы вызываете \t{exit(0)}, умирает весь процесс со всеми потоками.
		\item Но в \t{main} можно сделать \t{pthread\_exit}, если очень хочется, тогда процесс не умрёт, пока живы потоки.
	\end{itemize}
\end{frame}

\begin{frame}[fragile]{Упражнение: сборка кода}
	\begin{enumerate}
		\item Качаем решение с \href{https://github.com/yeputons/fall-2016-paradigms/raw/master/161019/sources/01-simple.c}{GitHub}.
		\item \texttt{gcc 01-simple.c -o 01-simple -std=gnu99 -Wall -Wextra -Werror -pthread} или аналог в вашей IDE.
		\item \texttt{./01-simple}
		\item Ожидаемый вывод:
\begin{verbatim}
Hello from thread! arg=1234
data is 1244
\end{verbatim}
	\end{enumerate}
\end{frame}

\begin{frame}{Несколько замечаний про C}
	\begin{itemize}
		\item На языке C лучше включить все предупреждения компилятора (warnings), в GCC это делают ключи \t{-Wall}, \t{-Wextra}.
		\item Если вы включили предупреждения "--- их лучше сразу трактовать как ошибки (\t{-Werror}), иначе быстро научитесь их игнорировать.
		\item Если аргумент функции не используется, после него в GCC надо писать \t{\_\_attribute\_\_((unused))}.
		\item Из функции всегда надо что-то вернуть (хотя бы \t{NULL}).
		\item \href{http://codeforces.com/blog/entry/17747}{Никогда} не начинайте название переменной с нижнего подчёркивания!
		\item Константы задаются при помощи \t{\#define SOME\_CONST (value)}.
		\item Встроенного типа \t{bool} нет, надо подключить \t{\#include~<stdbool.h>}
	\end{itemize}
\end{frame}

\begin{frame}[t]{Иллюстрация}
	\begin{center}
		\includegraphics[scale=0.5]{no-werror-meme.jpg}
	\end{center}
\end{frame}

\begin{frame}[t]{Несколько замечаний про pthread}
	Про потоки и pthread:
	\begin{itemize}
		\item Использовать \t{void* arg} и возвращаемое значение для передачи данных необязательно.
		\item Вся память внутри процесса одинаково доступна всем потокам на чтение и запись.
		\item \t{void* arg} возникает только тогда, когда надо запустить потоки на разных данных.
		\item Что произойдёт, если мы забудем \t{join} и \t{main} завершится до начала \t{worker}?
			\only<2->{Неопределённое поведение "--- \t{worker} попытается изменить переменную \t{data}, которая уже исчезла.}
	\end{itemize}
\end{frame}

\begin{frame}{Кто освобождает ресурсы?}
	На самом деле в pthread есть два типа потоков: joinable и detached.

	Joinable:
	\begin{itemize}
		\item Тип по умолчанию.
		\item На таком потоке должен быть ровно один раз вызыван метод \t{pthread\_join}, который освободит ресурсы и сообщит, что поток вернул.
		\item Если не вызвать "--- ресурсы не будут освобождены до конца программы.
		\item Если вызвать дважды "--- второй вызов может уронить программу или вернуть неверный результат.
	\end{itemize}

	Detached:
	\begin{itemize}
		\item Система автоматически освободит ресурсы как только поток завершится.
		\item Нельзя вызывать \t{pthread\_join} и получать возвращаемое значение "--- его негде хранить после окончания работы.
	\end{itemize}
\end{frame}

\begin{frame}{В других системах}
	\begin{itemize}
		\item Joinable/detached также используется в Java.
		\item В Windows (не в pthread под Windows!) другая концепция:
			\begin{itemize}
				\item Указатель на поток "--- сложный объект, который надо запрашивать у ОС и освобождать (как \t{FILE*}), а не просто переменная.
				\item Ресурсы потока освобождаются, когда он завершился и на него больше нет указателей.
				\item Нет разделения joinable/detached.
				\item Если кто-то может спросить состояние потока "--- у него есть указатель, значит, ресурсы потока ещё не освобождены.
			\end{itemize}
	\end{itemize}
\end{frame}

\begin{frame}{Упражнение}
	\begin{enumerate}
		\item Измените код так, чтобы \t{data} стала глобальной переменной (после этого \t{arg} не нужен).
		\item Вызовите \t{pthread\_detach} на втором потоке после запуска.
		\item Убедитесь, что программа упала.
		\item Уберите вызов \t{pthread\_join} и \t{printf} из основного потока.
		\item Убедитесь, что второй поток не всегда успевает отработать.
		\item Добавьте вызов \t{pthread\_exit} в \t{main}.
		\item Убедитесь, что после приложение перестало закрываться до окончания работы всех потоков.
	\end{enumerate}
	\href{https://github.com/yeputons/fall-2016-paradigms/raw/master/161019/sources/02-detached.c}{Код}
\end{frame}
